%% Annexe C : Modèles et procédures
\chapter{Modèles et Procédures}
\label{ann:templates}

\section*{Consigne de première intervention}
\begin{enumerate}
    \item Ne pas éteindre, ne pas redémarrer, ne pas « nettoyer ».
    \item Isoler la machine du réseau (sans l'éteindre).
    \item Prévenir le responsable de la réponse à incident.
    \item Noter l'heure et ce qui a été constaté.
    \item Attendre la capture des preuves avant toute action.
\end{enumerate}

\section*{Fiche de constat}
\begin{tableaucours}{L{3.4cm}L{8.8cm}}{%
    \tableauHeader{Rubrique} & \tableauHeader{Contenu}
}
Identifiant & T-xx, titre court \\
Constatation & Le fait vérifiable \\
Preuve & Artefact, condensat, source \\
Analyse & Interprétation, niveau de confiance \\
Gravité & Cotation justifiée \\
Recommandation & Correction concrète \\
\end{tableaucours}

\section*{Fiche de chaîne de custody}
\begin{tableaucours}{L{3.4cm}L{8.8cm}}{%
    \tableauHeader{Rubrique} & \tableauHeader{Contenu}
}
Pièce & Numéro, description, condensat \\
Acquisition & Qui, quand, comment, outil \\
Transferts & De qui à qui, quand, pourquoi \\
Conservation & Lieu, protection \\
\end{tableaucours}

\section*{Plan du rapport d'incident}
\begin{enumerate}
    \item Synthèse pour la direction.
    \item Chronologie des faits.
    \item Analyse (constatation / interprétation).
    \item Impact et obligations (ANTIC, données).
    \item Actions~: confinement, éradication, rétablissement.
    \item Retour d'expérience et recommandations.
\end{enumerate}
